\section{Práctica de ingeniería: cómo diseñar un sistema de test-time scaling}

\subsection{Primero hay que determinar si la tarea se presta a la expansión}

\begin{enumerate}
    \item ¿El modelo base es capaz de producir una respuesta correcta al menos de forma ocasional, es decir, $p>0$?
    \item ¿Existe un verifier más barato y más confiable que la generación?
    \item ¿El valor de la solicitud es suficiente para cubrir los tokens, la latencia y la energía adicionales?
    \item ¿Los candidatos pueden generarse en paralelo, o es necesario leer el estado del entorno de forma secuencial?
    \item ¿Las acciones fallidas son reversibles y se pueden reintentar de forma segura?
\end{enumerate}

\subsection{La regla de paro es más importante que el presupuesto máximo}

El sistema puede detenerse anticipadamente bajo las siguientes condiciones:

\begin{itemize}
    \item la puntuación del verifier supera un umbral confiable;
    \item varios verifiers independientes llegan a un consenso;
    \item la mejora marginal de calidad de un nuevo candidato es menor que su costo;
    \item revisiones sucesivas no corrigen el error clave;
    \item se alcanza el límite de latencia, costo o seguridad.
\end{itemize}

\begin{importantbox}{Controlador de presupuesto en línea}
Un sistema ideal no usa un best-of-64 fijo, sino que primero utiliza un presupuesto pequeño para estimar la dificultad del problema y la distribución de los candidatos, y después decide si conviene expandir. Los problemas sencillos se resuelven rápido, los problemas difíciles pero verificables reciben más cómputo, y los problemas que el modelo simplemente no puede resolver se escalan a tiempo hacia un modelo más fuerte o hacia intervención humana.
\end{importantbox}

\subsection{Requisitos de observabilidad}

Se debe registrar, para cada candidato, el modelo que lo generó, la semilla aleatoria, el número de tokens, la puntuación de verificación, la justificación de la selección, los resultados de las herramientas y el motivo del paro. De lo contrario, solo se observa la respuesta final, sin poder distinguir entre fallas de generación, fallas de verificación, fallas de la estrategia de búsqueda y presupuesto insuficiente.

\subsection{Resumen de la sección}

Test-time scaling es un problema de sistemas de control. Una implementación entregable requiere definir con claridad el espacio de candidatos, la señal de verificación, la función de presupuesto y la regla de paro, y evaluarse en conjunto mediante pass@1 de extremo a extremo, costo y latencia.
